\section{Interaction Detection}
\label{sec:od_physics}
\par
An energy deposit in the GdLS will produce a light yield that can be described by Birk's Law \cite{birks_law_ref};
\begin{equation} 
    \frac{dL}{dx} = S \frac{\frac{dE}{dx}}{1 + k_{B}\frac{dE}{dx} + C\frac{dE}{dx}^2}
    \label{eq:birkslaw}
\end{equation}
Here, L is the the light yield, S is the scintillation efficiency, $\frac{dE}{dx}$ is the energy loss of the particle per path length, and $k_{B}$ and $C$ are Birk's law parameters which differ depending upon the source; $\gamma$'s, electron recoils, etc...
These parameters have been well studied for cocktails of similar makeup to the LZ GdLS with identical PPO and bis-MSB concentrations \cite{ls_alpha_quenching_ref,ls_proton_quenching_ref}.
The low amount of Gadolinium-doping means that these should not differ significantly from that.
The parameters for GdLS are shown in table \ref{tab:Birks_law_parameters}.
\begin{table}[!htbp]
    \centering
    \begin{tabular}{c | c | c | c }
                   & S (photons/MeV) & $k_{B}$ & $C$ \\ \hline
    $\alpha$       &                 & 4.63E-3 & 1.77E-6 \\
    $\beta/\gamma$ & 9,000           & 0.03    & 0 \\ 
    proton         &                 & 8.26E-3 & 0
    \end{tabular}
    \caption{GdLS parameters for Birk's law for the LZ OD.}
    \label{tab:Birks_law_parameters}
\end{table} 

\par
The source of the interactions will generally be quenched, with $\alpha$'s producing the least amount of light per energy deposited.
This is summarised in Figure \ref{fig:gdls_quenching}.

\begin{figure}[!htbp]%
    \centering
    \includegraphics[height=8cm]{Figures/Placeholder.png}
    \caption{Quenching}
    \label{fig:gdls_quenching}
\end{figure}

\subsection{Neutron Propagation}
\par
In the lab-frame, the ratio between a neutrons initial energy ($E$) and final energy ($E'$) when colliding with an at rest nucleus of mass A follows;
\begin{equation}
    \frac{E'}{E} = \frac{A^2 + 1 + 2A\cos{\theta}}{(A + 1)^2}
\end{equation}
where $\theta$ is the scattering angle in the CoM.
Now, if the nucleus is a hard sphere, so independent of $\theta$, then an elastic scatter follows
\begin{equation}
    E^{'} = E(\frac{A-1}{1})^{2}
\end{equation}
If A = 1 (so is a proton), then it reduces further to $E^{'} = \frac{1}{2}E$, meaning that the neutron loses half its energy on average.
The scattered proton has energy $E - E^{'}$ and is quenched.
Neutrons predominately collide effectively on Hydrogen atoms \cite{neutron_thermalisation_and_capture_ref}.

\par
Once thermalised, a neutron a neutron will capture most often on a Gadolinium isotope (87\% of the time), and a Hydrogen on the remainder\footnote{Other isotopes can capture the neutron but the probability is sufficiently low that they are not included here.}.
Despite only making up a small fraction of the GdLS cocktail, the capture cross-section of Gadolinium is orders of magnitude above that of Hydrogen to suffice, as can be seen in Figure \ref{fig:gdls_capture_cross_sections}.

\begin{figure}
	\begin{tikzpicture}
		\centering
		\begin{axis}[%point meta max=150,
			%point meta min=0.0,
			view={0}{90},
			ylabel={Cross-Section (barns)},
			xlabel={Incident Energy (MeV)},
			width=15cm,
			height=8cm,
			ymode=log,
			xmode=log,
			]
			\addplot[blue] table[x=energy,y=barns] {Data/GdLS_Physics/capture_cross_sections/h.dat};
			\addplot[orange] table[x=energy,y=barns] {Data/GdLS_Physics/capture_cross_sections/gd155.dat};
			\addplot[green] table[x=energy,y=barns] {Data/GdLS_Physics/capture_cross_sections/gd157.dat};
			\legend{H, ${}^{155}$Gd, ${}^{157}$Gd};
		\end{axis}
	\end{tikzpicture}
\caption{Capture cross sections for the most significant capture particles in GdLS cocktail. 
A neutron has an increase capture probability if it's energy is close to that of the isotope energy. 
Values from ENDF/B-VIII.0 \cite{endf_bv3_ref}.}
\label{fig:gdls_capture_cross_sections}
\end{figure}
		

\par
Once a neutron is captured, the new isotope de-excites as either direct $\gamma$ emission or internal-conversion resulting in ejected auger electrons or characteristic x-rays.
For a n-H capture, a single 2.2MeV $\gamma$ is emitted, which may leave the GdLS without scattering, and so the neutron is not detected.
However, n-Gd capture produces $\approx$ 8MeV worth of particles, typically as a number of $\gamma$'s.
Between the excited and ground states of the Gadolinium isotopes ${}^{156}$Gd and ${}^{158}$Gd there are in excess of $10^{6}$ levels.
When the simulation framework was setup for the LZ experiment, GEANT4 was not adequate for modelling high atomic mass de-excitation's where there are a huge number of possible cascade combinations \cite{ucsb_gdls_dicebox_simulations_ref}, and it has not been a problem for other experiments which used Gadolinium-doped LS, such as DayaBay \cite{dayabay_overview_ref} and RENO \cite{reno_overview_ref}, as the quantity of scintillator is large enough that the entirety of the de-excitation energy is generally deposited in the medium.
However, the LZ OD is significantly thinner than those, with a non-insignificant probably of not all of the energy being deposited in the GdLS.
As such LZ incorporated DICEBOX \cite{dicebox_simulations_ref} simulations, which are sampled allowing for a more realistic simulation \cite{lz_simulations_ref}.
A detailed comparison of the DICEBOX result against GEANT4 can be found in \cite{ucsb_gdls_dicebox_simulations_ref}.


\par
Using DICEBOX, the proportion of how the energy is dissipated in the isotope is shown in Figure \ref{fig:od_gdls_dicebox_deexcitation}, and the resultant $\gamma$ spectrum in Figure \ref{fig:gd_capture_resulting_gamma_spectrum}.
For both ${}^{156}Gd$ and ${}^{158}Gd$ the de-excitation occurs in an average of 4-5 steps, with the vast majority of the energy going directly to $\gamma$-emission.
Because of this, the de-excitation of Gadolinium is often quoted as "4.7 $\gamma$'s".
This highlights the importance of doping the scintillator as there is a much higher probability of being able to observe a n-Gd capture with multiple $\gamma$'s than it is a n-H capture with a single $\gamma$.

\begin{figure}[!htbp]%
\centering
\begin{tikzpicture}
\centering
  \begin{groupplot}[%view={0}{90},
    group style = {group size = 2 by 1}]
    
    \nextgroupplot[
            xlabel=N. Steps,
            ylabel=,
            width=0.45\textwidth,
            height=6cm,
            xmin=0,
            xmax=13,
            ymin=0,
            legend pos=north east,]
            \addplot[red, const plot]
                    table [x=Lower,y=Weight]
                    {Data/GdLS_Physics/DICEBOX/gd156_n_steps.txt};
                %\addlegendentry{$Gd^{156}$};
            \addplot[blue, const plot]
                    table [x=Lower,y=Weight]
                    {Data/GdLS_Physics/DICEBOX/gd158_n_steps.txt};
                %\addlegendentry{$Gd^{158}$};
    
    \nextgroupplot[
            xlabel=Energy Transfer location,
            ylabel=Percentage of steps,
            width=0.45\textwidth,
            height=6cm,
            xtick=data,
            xticklabels from table={Data/GdLS_Physics/DICEBOX/gd158_energy_conversion_to_shells.txt}{Label}]
            \addplot[red, const plot]
                    table [x=Bins, y=Weight]
                    {Data/GdLS_Physics/DICEBOX/gd158_energy_conversion_to_shells.txt};
                %\addlegendentry{$Gd^{156}$};
            \addplot[blue, const plot]
                    table [x=Bins, y=Weight]
                    {Data/GdLS_Physics/DICEBOX/gd156_energy_conversion_to_shells.txt};
                %\addlegendentry{$Gd^{158}$};
        
  \end{groupplot}
\end{tikzpicture}
\caption{Empty caption}
\label{fig:od_gdls_dicebox_deexcitation}
\end{figure}

\begin{figure}[]
    \centering
    \begin{tikzpicture}
        \begin{axis}[
            %title=$\gamma$ energy spectrum from $Gd^{155}$ and $Gd^{157}$ neutron captures,
            xlabel=Energy (KeV),
            ylabel=Emission (Arb.),
            width=15cm,
            height=6cm,
            xmin=0,
            xmax=8000,
            ymin=0,
            legend pos=north east,]
            \addplot[red, const plot]
                    table [x=Lower,y=Weight]
                    {Data/GdLS_Physics/DICEBOX/gd156_gamma_spectrum.txt};
                \addlegendentry{${}^{156}Gd$};
            \addplot[blue, const plot]
                    table [x=Lower,y=Weight]
                    {Data/GdLS_Physics/DICEBOX/gd158_gamma_spectrum.txt};
                \addlegendentry{${}^{158}Gd$};
        \end{axis}
    \end{tikzpicture}
    \caption{DICEBOX simulation of the energy released by $\gamma$s from de-excitation after neutron capture for ${}^{156}Gd$ and ${}^{158}Gd$}
    \label{fig:gd_capture_resulting_gamma_spectrum}
\end{figure}